% Financial credit-risk bootstrap experiment for WonderFullyHE.
% This file is intentionally self-contained so it can be copied into a thesis/report
% and extended with measured tables and plots later.
\section{Эксперимент: конфиденциальная оценка ожидаемых кредитных потерь}

\subsection{Цель эксперимента}

Целью эксперимента является демонстрация применения гомоморфного шифрования
CKKS для вычисления финансового показателя ожидаемых кредитных потерь
(\emph{Expected Loss}, EL) без раскрытия индивидуальных значений вероятности
дефолта. Дополнительно в вычислительную цепочку включается процедура
бутстраппинга, после которой вычисления над шифртекстом продолжаются.

Для $i$-го кредита ожидаемый убыток задаётся выражением
\[
    EL_i = PD_i \cdot EAD_i \cdot LGD_i \cdot s_i,
\]
где $PD_i$ --- вероятность дефолта,
$EAD_i$ --- величина требования на момент дефолта,
$LGD_i$ --- доля потерь при дефолте,
а $s_i$ --- коэффициент стресс-сценария.
Агрегированный показатель портфеля равен
\[
    EL_{\mathrm{portfolio}}
    = \sum_{i=1}^{n} EL_i.
\]

\subsection{Связь с моделью кредитного скоринга}

В полном сценарии вероятность дефолта может быть получена из зашифрованных
признаков заёмщика. Сначала вычисляется линейный скоринговый показатель
\[
    z_i = w_0 + \sum_{j=1}^{m} w_j x_{ij},
\]
после чего нелинейная функция активации заменяется полиномиальной
аппроксимацией
\[
    PD_i \approx p(z_i)
    = a_0 + a_1 z_i + a_2 z_i^2 + \ldots + a_d z_i^d.
\]
Именно последовательные умножения при вычислении $p(z)$ расходуют уровни
modulus chain CKKS. После уменьшения доступной вычислительной глубины
бутстраппинг должен восстановить возможность продолжения вычислений.

Текущая реализация эксперимента фиксирует границу непосредственно перед
бутстраппингом: в шифртексте уже находятся значения $PD_i$. Это сделано
намеренно, поскольку production-путь библиотеки на текущем этапе не имеет
полного доказанного lift/security certificate для произвольного обычного
шифртекста. Таким образом, эксперимент проверяет именно связку
\[
    \operatorname{Enc}(PD)
    \longrightarrow \operatorname{Bootstrap}
    \longrightarrow \operatorname{EL}
    \longrightarrow \operatorname{Aggregate}.
\]
После устранения production-блокеров перед этой границей можно добавить
гомоморфное вычисление $z_i$ и полиномиальной аппроксимации $p(z_i)$ без
изменения финансовой части эксперимента.

\subsection{Экспериментальная схема}

Используется пакетная обработка восьми кредитов. Вероятности дефолта
размещаются в CKKS-слотах:
\[
PD =
(0.004,\ 0.008,\ 0.012,\ 0.016,\
 0.020,\ 0.024,\ 0.028,\ 0.032).
\]
Публичные коэффициенты $EAD_i LGD_i s_i$ задаются вектором
\[
L =
(42000,\ 56000,\ 67500,\ 72000,\
 82500,\ 90000,\ 104000,\ 117000).
\]
После бутстраппинга выполняется покомпонентное умножение
\[
    C_{EL}
    = C_{PD}^{\mathrm{refresh}} \odot L,
\]
после чего значения всех активных слотов суммируются при помощи ротаций и
сложений:
\[
    C_{\mathrm{portfolio}}
    =
    C_{EL}
    + \operatorname{Rot}(C_{EL},1)
    + \operatorname{Rot}(C_{EL},2)
    + \operatorname{Rot}(C_{EL},4).
\]
После расшифрования первый активный слот содержит приближение
$EL_{\mathrm{portfolio}}$.

\subsection{Параметры исследовательского стенда}

Для запуска полного исполняемого пути бутстраппинга используется существующий
в проекте test-only стенд:
\[
    N=16,\qquad
    \Delta_0 = 2^{49},
\]
с пятнадцатью 50-битными простыми модулями. Стенд использует
\texttt{security level = none} и синтетический шифртекст вида $c=(m,1)$.
Для lift bound явно задаётся предположение $K=1$ типа
\texttt{TestFixtureAssumption}.

Следовательно, данный эксперимент является демонстрацией корректности
вычислительного pipeline и влияния бутстраппинга на дальнейшие финансовые
вычисления, но не является доказательством production-безопасности
параметров. Результат глобальной сертификации не должен интерпретироваться
как \texttt{Certified} для промышленного применения.

\subsection{Реализация}

Эксперимент реализован в файле
\begin{verbatim}
apps/experiment_credit_risk_bootstrap.cpp
\end{verbatim}
и собирается только при одновременно включённых
\texttt{BUILD\_TESTING} и \texttt{M2424\_ENABLE\_EVALMOD\_ANALYSIS}.

Пример сборки:
\begin{verbatim}
cmake -S . -B build-risk \
  -DCMAKE_BUILD_TYPE=Release \
  -DBUILD_TESTING=ON \
  -DM2424_ENABLE_EVALMOD_ANALYSIS=ON

cmake --build build-risk -j
./build-risk/bin/experiment_credit_risk_bootstrap
\end{verbatim}

Программа выводит:
\begin{itemize}
    \item глобальный статус и первый gate бутстраппинга;
    \item аналитическую границу ошибки $E_{\mathrm{boot}}$;
    \item наблюдаемую ошибку восстановленных $PD_i$;
    \item chain index сразу после бутстраппинга и после финансового этапа;
    \item эталонный plaintext Expected Loss;
    \item результат вычисления над шифртекстом;
    \item абсолютную и относительную ошибку.
\end{itemize}

\subsection{Метрики}

В экспериментальную таблицу рекомендуется включить следующие показатели:
\[
\begin{array}{|l|l|}
\hline
\text{Показатель} & \text{Назначение}\\
\hline
E_{\mathrm{boot}} &
\text{сертифицированная верхняя граница ошибки бутстраппинга}\\
\hline
E_{\mathrm{obs}}^{PD} &
\max_i |\widehat{PD}_i-PD_i|\\
\hline
E_{\mathrm{abs}}^{EL} &
|\widehat{EL}_{\mathrm{portfolio}}-EL_{\mathrm{portfolio}}|\\
\hline
E_{\mathrm{rel}}^{EL} &
E_{\mathrm{abs}}^{EL}/|EL_{\mathrm{portfolio}}|\\
\hline
L_{\mathrm{after\ boot}} &
\text{уровень modulus chain сразу после refresh}\\
\hline
L_{\mathrm{final}} &
\text{уровень после расчёта Expected Loss}\\
\hline
T_{\mathrm{boot}} &
\text{время выполнения бутстраппинга}\\
\hline
T_{\mathrm{EL}} &
\text{время пост-бутстраппингового финансового расчёта}\\
\hline
\end{array}
\]
В текущей версии программы временные метрики можно добавить отдельными
замерами \texttt{std::chrono} вокруг \texttt{bootstrap.apply} и финансового
этапа.

\subsection{Ожидаемая интерпретация результата}

Эксперимент должен показывать три свойства системы.

Во-первых, значения вероятности дефолта не раскрываются серверу во время
вычисления. Сервер выполняет бутстраппинг и последующий расчёт Expected Loss
над шифртекстами.

Во-вторых, после бутстраппинга вычисления продолжаются над обновлённым
шифртекстом: выполняются plaintext-vector multiplication, rescale, ротации и
суммирование.

В-третьих, получаемый результат можно непосредственно сопоставить с
plaintext-эталоном:
\[
    EL_{\mathrm{portfolio}}^{\mathrm{plain}}
    \quad\text{и}\quad
    EL_{\mathrm{portfolio}}^{\mathrm{CKKS}}.
\]
Это позволяет одновременно оценивать численную ошибку, расход уровней и
стоимость применения бутстраппинга.

\subsection{Дальнейшее расширение}

После завершения production-пути бутстраппинга эксперимент можно расширить до
полной цепочки
\[
\operatorname{Enc}(x_i)
\to z_i
\to p(z_i)
\to \operatorname{Bootstrap}
\to EL_i
\to EL_{\mathrm{portfolio}},
\]
где $p(z)$ является полиномиальной аппроксимацией sigmoid. Тогда необходимость
бутстраппинга будет возникать естественно из-за мультипликативной глубины
скоринговой модели, а не задаваться только как экспериментальная граница
между двумя частями вычисления.
